\chapter*{TÓM TẮT}
\addcontentsline{toc}{chapter}{Tóm tắt}

\vspace{0.5cm}

\noindent\textbf{Từ khóa:} Game Balancing, Thuật toán di truyền, NSGA-II, Roguelike Deckbuilder, Tối ưu hóa đa mục tiêu, Procedural Content Generation

\vspace{1cm}

Cân bằng game (game balancing) là một thách thức lớn trong phát triển trò chơi điện tử, đặc biệt với thể loại Roguelike Deckbuilder nơi mỗi lượt chơi là độc nhất với các yếu tố ngẫu nhiên. Phương pháp truyền thống dựa vào kinh nghiệm và thử nghiệm thủ công tốn nhiều thời gian và dễ bỏ sót các kịch bản quan trọng.

Đề tài này nghiên cứu và phát triển một hệ thống tự động cân bằng game sử dụng thuật toán di truyền (Genetic Algorithm - GA). Hệ thống bao gồm: (1) một Roguelike Deckbuilder prototype hoàn chỉnh với hệ thống chiến đấu, quản lý bộ bài, sinh bản đồ và hệ thống di vật; (2) AI agent tự động chơi game để thu thập dữ liệu; (3) hai phương pháp tối ưu hóa: GA đơn mục tiêu với GA truyền thống tối ưu \~{}500 tham số và NSGA-II đa mục tiêu tối ưu \~{}50 tham số phân cấp.

Thực nghiệm với 30 thế hệ và 100 cá thể mỗi thế hệ cho thấy phương pháp NSGA-II đa mục tiêu vượt trội hơn: tiết kiệm 33\% thời gian tính toán (49 phút so với 73 phút), tăng Balance Score từ 0.754 lên 0.960 (+27.3\%), Engagement Score từ 0.655 lên 0.900 (+37.4\%), và đạt Coherence score tối đa 1.000 (+17.2\%). Phương pháp này tìm được Pareto front gồm nhiều giải pháp không bị chi phối, cho phép người dùng chọn sự đánh đổi phù hợp giữa Balance, Engagement và Coherence. Số lượng thẻ bài khả dụng tăng 16.7\% (12 lên 14) và Build Variety tăng 34.8\% (0.705 lên 0.950).

Hệ thống chứng minh hiệu quả của việc kết hợp kiến thức chuyên môn (gom nhóm tham số phân cấp) với tối ưu hóa đa mục tiêu trong cân bằng game. Công cụ có tiềm năng lớn để tự động hóa quá trình cân bằng, tiết kiệm hàng trăm giờ điều chỉnh thủ công cho các nhà phát triển game.
